% SPDX-FileCopyrightText: 2026 Will Cook
% SPDX-License-Identifier: CC-BY-4.0
\section{Explain the choices in a proof}\label{sec:revisions}
\papersectiontarget{exposition-revisions}

These cases concern attainment and the accuracy an inference requires.
They were accepted in the short-paper R6 round of 30 September 2026, not the
earlier September series with the same round numbers. The records retain
the passages at commit \texttt{18cedaddedd1}, its full identity and source
digests~\cite{ledger}.

\subsection{From repeated choice to an attained infinite sum}\label{sec:remainder}

The \#251 paper contains a construction in which a target $x$ is represented
by chosen digits $d_j$ and weights $w_j$. A finite remainder $\rho_j$ is kept
between $0$ and a bound $F_j$. The old passage began:
\begin{quote}
Given $x\in[0,F_0]$, we choose $d_0$ so that
$x-d_0w_0\in[0,F_1]$ and repeat this choice for each successive remainder.
\end{quote}
This leaves two questions: why is the next digit available, and why does
continuing the choices represent the whole target? The reviewed addition
answers them separately:
\begin{quote}
The covering just proved makes this possible at every stage. The invariant
$0\le\rho_j\le F_j$ and $F_j\to0$ then give
\end{quote}
and continues with the limiting identity. The first sentence says why the
construction can continue. The second says why its partial sums reach the
chosen target: the part still missing is forced to zero. The record is
\caseid{r6-251-style\_rules-04}; its full identifier includes the series date.

Crmari\'c and Kova\v{c}'s proof of Lemma~4(a), in the first version of their
paper on sums of reciprocals, provides a close model~\cite[pp.~4--5]{crmaric}.
They keep a finite remainder in an interval and use convergence to obtain
the represented value. What transfers is the separation of two questions:
can the next choice be made, and does the resulting sequence attain the target?

The \#251 construction must justify its covering, invariant and $F_j\to0$,
as well as sparsity and arithmetic constraints. Neither a finite digit
example nor the literary comparison supplies those proofs. The integration
preserved these obligations and the unresolved status of the original
prime-gap problem.

\subsection{State how accurate the estimate must be}\label{sec:estimate}

In the \#269 argument, a rationality assumption places normalized tails
$Y_a$ in the finite set
\[
             (K^{-1}\mathbb Z)\cap(0,1).
\]
Here $K\ge1$ is the integer denominator in the rationality assumption.
Finiteness forces two tails to agree; it does not force their successors to
agree. The accepted revision makes the remaining task explicit: recover the
next block and tail from the present tail. Only then does equality propagate.

Han\v{c}l and Tijdeman begin their proof of Theorem~2.1 with an integral
scaled remainder obtained from rationality~\cite[p.~373]{hancl}. The \#269
revision likewise identifies the arithmetic object before estimating it;
recovery is its own additional task, not supplied by the cited proof.

The next difficulty is quantitative. The maps
$G_{53}(t)=(9+t)/30$ and $G_5(t)=(3+t)/10$ encode two possible blocks of
prime-power jumps. Their images overlap on the initial tail range $[0,1]$, so a value in the
overlap does not identify its block. Both maps are increasing: to separate
their images, compare the upper end of one with the lower end of the other.
With the lower bound $3/10$ already established, the revision states the
required accuracy first:
\begin{quote}
With that lower endpoint, separating these two images requires an upper
endpoint $U$ satisfying
\[
 G_{53}(U)<G_5(3/10),\qquad
 \frac{9+U}{30}<\frac{33}{100},\qquad U<\frac9{10}.
\]
This specifies the improvement needed from the arithmetic sequence.
\end{quote}
The case is \caseid{r6-269-style\_rules-03}. The threshold is sufficient for
this argument, not necessary for every possible method. The return's Knuth
analogy concerns presentation; the separation inequality supplies the mathematics.

The source obtains the required improvement from the spacing of powers of $5$.
At a block containing such a power the tail is at most $7/15$. Every three
consecutive blocks contain at least one such block; moving backwards at most
twice, each step is bounded by $t\mapsto(1+t)/2$. This gives $U=13/15<9/10$, for the actual
prime-power sequence, not arbitrary block words. The source separates all five
block images on this interval. The image identifies the block; its affine
inverse determines the next tail. Equality can now propagate~\cite[proof of the distinct-height theorem]{paper269}.

\subsection{An example must reveal the relevant difference}\label{sec:multiplicity}

The same manuscript distinguishes summing over smooth integers from summing
over the distinct values of their running least common multiple. Here the
allowed prime factors are $2,3,5$; a positive integer with no other prime
factor is called $5$-smooth. For $x\ge1$ its running least common multiple is
\[
 H(x)=\operatorname{lcm}\{u\le x:u\text{ is }5\text{-smooth}\}
     =2^{\lfloor\log_2x\rfloor}
      3^{\lfloor\log_3x\rfloor}
      5^{\lfloor\log_5x\rfloor}.
\]
The formula selects the largest power of each allowed prime not exceeding
$x$~\cite[definition of running least common multiples]{paper269}. For example,
$H(16)=16\cdot9\cdot5=720$. These are the quantities called heights in the
revision. The seven $5$-smooth integers in $[16,32)$ give:
\begin{center}
\begin{tabular}{@{}lll@{}}
\toprule
$5$-smooth integers $u$ & $H(u)$ & Sum of $1/H(u)$\\
\midrule
$16,18,20,24$ & $720$ & $4/720$\\
$25$ & $3600$ & $1/3600$\\
$27,30$ & $10800$ & $2/10800$\\
\bottomrule
\end{tabular}
\end{center}
The map from a smooth integer to its height is not one-to-one. Summing over
smooth integers counts every visit to a height; summing over distinct heights
counts it once. The multiplicities $4$ and $2$ are the sizes of the corresponding
fibres. They explain why the sums cannot be identified merely because the
same heights occur in both.
The record \caseid{r6-269-style\_rules-04} also retains an endpoint warning:
the distinct-height argument uses $(16,32]$, whereas the table uses
$[16,32)$. Matching the endpoint conventions is a separate task from
preserving multiplicity.

The repeated heights make this table useful: they expose the counting
convention. Halmos's concrete cases and Gowers's examples-first discussion
suggest that choice~\cite[\S4]{halmos}\cite{gowers}; they prescribe no order
for every reader. The infinite-series theorem still needs its proof.


\section{From a formula to the sentence that explains it}\label{sec:worked-explanations}

These cases answer a reader's question from an existing formula. They come
from eight additional editorial returns (\#1049 is labelled R9), checked
against current manuscript sources. They add no theorem, formal coverage
or release status.

\subsection{How does the target determine the finite test?}\label{sec:certificate}

A certificate is easier to understand when the reader sees the set it must
certify before seeing its arithmetic form. In the \#251 paper, two tail
differences satisfy
\[
 D'=2D-\delta,\qquad \delta=2s,\qquad s\in\{-1,1\}.
\]
The already proved local criterion asks for $1/2<sD<1$; it then gives
$-1<sD'<0$. The finite calculation supplies integers $A,Q,B$ with $Q>0$,
$B\ge0$ and $|QD-A|\le B$~\cite[section on finite separation]{paper251}.
It is tempting to write only that two ensuing inequalities certify the target.
The missing explanation is the enclosure itself:
\[
 sD\in\left[\frac{sA-B}{Q},\frac{sA+B}{Q}\right].
\]
Multiplication by either sign preserves the error bound. The lower endpoint
must exceed $1/2$, and the upper endpoint must be below $1$. Thus the entire
closed enclosure lies in $(1/2,1)$ exactly when
\[
             2sA-Q>2B,\qquad Q-sA>B.
\]
The two inequalities express the two endpoint gaps. The factor $2$ in the
first comes from the lower endpoint $1/2$, not an arbitrary safety margin.
Strictness matters because the boundary values give $sD'=-1$ or $0$.

Now separate a failed enclosure from a failed claim. Take $D=3/4$, $s=1$,
$Q=4$, $A=3$ and $B=1$. The error bound is valid, but the enclosure
$[1/2,1]$ touches both forbidden endpoints. Neither strict test passes,
although $D$ lies in the target interval. With these same $D,Q,A$, the sharper
bound $B=0$ would pass both tests. The failed enclosure therefore need not
mean a failed claim. The enclosure even has width less than one and still
contains an integer: its position, not just its width, matters.

The \#68 remainder has a related, one-sided geometry. It is
$A_N+M(S-H_N)$, where $S$ is the series sum, $H_N$ its $N$th partial sum,
and $M>0$ scales the strictly positive omitted tail. To put this remainder
between consecutive integers, it suffices to prove
\[
             0<M(S-H_N)<\lfloor A_N\rfloor+1-A_N.
\]
The left inequality puts the remainder above $\lfloor A_N\rfloor$; the right
puts it below $\lfloor A_N\rfloor+1$~\cite[section on nonintegrality]{paper68}.
A tail smaller than one could still cross the next integer; it must be
smaller than this particular gap. At an integral $A_N$ the gap is $1$. Replacing the strict successor by
$\lceil A_N\rceil$ would make it zero. The endpoint distinguishes two
superficially interchangeable conventions.

\subsection{Why was this base or kernel chosen?}\label{sec:formula-choice}

A change of notation should reveal a constraint or save an operation. In one
\#249 comparison construction, only the even-indexed coefficients may change.
Writing $e_n$ for the coefficient changes, their contribution to a binary series is
\[
                 \sum_{m\ge1} e_{2m}\,2^{-2m}
                 =\sum_{m\ge1}e_{2m}\,4^{-m}.
\]
The permitted support has selected base four. Stating this before the digit
construction explains the choice and shows immediately why the odd coefficients
remain untouched~\cite[comparison-sequence construction]{paper249long}.
The identity alone says nothing about which corrections can be attained under
additional coefficient bounds; those are separate assertions of the construction.

The \#257 divisor kernel admits an equally direct introduction. Fix integers
$N\ge0$ and $d\ge1$, and let $B>1$. If $r\ge1$ satisfies $d\mid N+r$,
then its possible values are
\[
       d-(N\bmod d),\quad 2d-(N\bmod d),\quad\ldots.
\]
Consequently
\[
 \sum_{\substack{r\ge1\\d\mid N+r}}B^{-r}
 =\frac{B^{N\bmod d}}{B^d-1}.
\]
This is the weight of the future positions at which $d$ divides the index.
It explains the kernel before averaging is introduced
\cite[finite means for shifted divisor tails]{paper257}. Divisibility applies
to $N+r$, not to $r$ alone. When $d\mid N$, the first positive offset is $d$,
so the same formula applies without adding an unwanted offset zero. In that paper $B=2^\alpha$, with $0<\alpha\le1$; estimates
must retain their dependence on $B-1$ as $\alpha$ becomes small.

In both examples the operation determines the notation. Even-indexed
binary weights become base-four weights; offsets to future indices divisible
by a fixed divisor form the progression giving the kernel. State that
reason before naming the formula. Tao's notation advice supports making
important dependencies visible and translating borrowed conventions~\cite{tao-notation};
the support and divisibility conditions supply the mathematical reasons here.

\subsection{What remains to justify after an index is fixed?}\label{sec:fixed-partition}

Fixing a summation index need not leave a fixed finite product. In the \#1049
moment-determinant argument, a summand is indexed by a partition
$\mu_1\ge\cdots\ge\mu_\ell>0$, with $\mu_k=0$ for $k>\ell$.
The base $q\in(0,1)$ and the partition are fixed while $N\to\infty$.
The ratios of shifted weights tend to one, but the Vandermonde quotient
still contains factors indexed by $k$ up to $N$
\cite[proof for geometric moment determinants]{paper1049}.

For $N\ge\ell$ and fixed $1\le j\le\ell$, cancel the common factors in
the part with $k>\ell$:
\[
 \prod_{k=\ell+1}^{N}
       \frac{1-q^{k-j+\mu_j}}{1-q^{k-j}}
 =\frac{\displaystyle\prod_{r=0}^{\mu_j-1}(1-q^{N+1-j+r})}
        {\displaystyle\prod_{r=0}^{\mu_j-1}(1-q^{\ell+1-j+r})}.
\]
The original product has $N-\ell$ factors, but each endpoint product has the
fixed length $\mu_j$. Its numerator tends to $1$; its denominator is fixed
and positive. The remaining factors have $j,k\le\ell$ and are independent
of $N$, so the whole quotient converges for this fixed partition. The
telescoping formula, proposed in the later return, justifies what factorwise
convergence in a growing product would not.

Passing the limit through the sum requires a second argument. Extend the
summand by zero when $\ell>N$, fixing the index set. The source's nonnegative
majorant is independent of $N$. Dropping the ordering of the positive parts
gives an upper bound by separate sums for each part. The $j$th sum is at most
a fixed constant times $q^{2j-1}$: geometric decay controls its polynomial
weight. Multiplying gives a bound for the total majorant at length $\ell$:
\[
                         K^\ell q^{\ell^2},
\]
where $K>0$ may depend on the fixed base and the fixed weight-bound constants.
Here $\ell^2=1+3+\cdots+(2\ell-1)$ is the minimum total displacement exponent
for $\ell$ positive parts. Their sizes have been summed out; their number
remains. The ratio of successive bounds is $Kq^{2\ell+1}\to0$, so the sum over
$\ell$ converges and dominated convergence applies.

The conclusions are distinct: the whole summand converges for a fixed partition;
a summable bound controls all partitions uniformly in $N$. Neither asserts
uniformity as $q\uparrow1$. The fixed base is part of the statement.

\subsection{Does a complete test provide the required witnesses?}\label{sec:detection}

In the \#249 test, fix a positive integer shift $h$ and an index $N\ge0$,
and write $D$ for that particular real tail difference. At each positive integer
depth $L$, the paper computes an integer $A_L$ and proves
\[
       |2^LD-A_L|\le B_L,\qquad B_L=N+h+L+2.
\]
An integer $D$ would put a multiple of $2^L$ in this closed enclosure. Hence
one can certify nonintegrality by checking
\[
       B_L<(A_L\bmod 2^L)<2^L-B_L,
\]
where the residue is chosen in $\{0,\ldots,2^L-1\}$
\cite[section on binary totient tails]{paper249}. The two strict margins
keep the enclosure away from the neighbouring multiples of $2^L$.

The test also detects every nonintegral fixed $D$ at some depth. Its gap
$\|D\|_{\mathbb R/\mathbb Z}$ from the nearest integer is then a positive
constant. The source proves that detection follows once
\[
                 2^L\|D\|_{\mathbb R/\mathbb Z}>2B_L.
\]
The factor $2$ pays for two distances: from $2^LD$ to the computed centre
$A_L$, then from that centre to an enclosure endpoint. Each is at most $B_L$;
the displayed gap keeps the whole enclosure away from multiples of $2^L$.
Since $2^L$ grows exponentially and $B_L$ linearly, a suitable depth exists.
Thus the test has a converse, not merely a sufficient condition.

The converse has two limits. An integral value never triggers the test;
failure at the depths tried does not certify integrality. Nor does detecting
a fixed nonintegral value supply the inputs needed for irrationality: for
every positive shift and every cutoff, some later index must have a nonintegral
difference. Existence of such an index is a separate question; increasing depth
only refines the test of a fixed index. Preserve the converse without mistaking
it for the missing existence statement.

Whenever a proof offers arbitrarily accurate certificates, state which object
stays fixed as the accuracy improves. If the object changes too, its distance
from the forbidden set may shrink, and a new comparison is needed. This is the
same discipline that made the fixed-base limit in
Section~\ref{sec:fixed-partition} readable.
